\chapter[4.3.a--b Indicadores, pruebas y resultados]{4.3.a--b Indicadores de desempeño, pruebas y análisis de resultados}
\label{ch:pruebas}

\section{4.3.a Indicadores de desempeño}

Los indicadores expresan el desempeño de los cuatro métodos en una misma escala cuantitativa \cite{martinRodellarControlAdaptativo}.

\subsection*{Indicadores Fase~1: Servoregulador}

\begin{longtable}{@{}P{0.22\linewidth}P{0.42\linewidth}P{0.24\linewidth}@{}}
\toprule
Indicador & Definición & Objetivo \\ \midrule
IAE servo & $\text{IAE} = \sum_{k=10}^{150}|r(k)-y(k)|$ & Minimizar.\\
Sobreimpulso $M_p$ & $M_p = 100\cdot\max(0,\,\max y - r)/r$\,\% & $< 5$\,\%.\\
Tiempo de estab. $t_s$ & Muestras hasta $|e| \leq 2\%$ de $r$. & Minimizar.\\
Tiempo de recuperación & Muestras para $|e| \leq 2.5\%$ tras fin perturbación. & Minimizar.\\
Energía servo & $E = \sum_{k=10}^{150}u(k)$ & No crecer sin mejora de IAE.\\
Variación de mando & $\sum_{k}|\Delta u(k)|$ & Minimizar (suavidad).\\
\bottomrule
\end{longtable}

\subsection*{Indicadores Fase~2: Coordinación aware R2ET}

\begin{longtable}{@{}P{0.24\linewidth}P{0.40\linewidth}P{0.24\linewidth}@{}}
\toprule
Indicador & Definición & Objetivo \\ \midrule
IAE aware & $\sum_{k}(|n_1-n^*|+|n_2-n^*|)$ & Minimizar.\\
IAE durante atasco & Mismo, solo en $k=95$--$118$. & Minimizar.\\
Ocupación máxima & $\max_{k,i}n_i(k)$ & $< 3$ (duro) y $< 2.7$ (preferible).\\
Margen de seguridad & $N_{max} - \max_{k,i}n_i(k)$ & Maximizar.\\
Tiempo en alarma & Muestras con $n_i > 2.7$. & Cero.\\
Violaciones & Muestras con $n_i > 3$. & Cero absoluto.\\
Productividad total & $\sum_{k}(y_{out,1}+y_{out,2})$ & Maximizar.\\
Energía aware & $\sum_{k}(u_1+u_2)$ & Equilibrar con productividad.\\
Desbalance medio & $\bar{b} = \overline{|n_1-n_2|}$ & $< 0.5$ piezas.\\
\bottomrule
\end{longtable}

\section{Escenarios de simulación}

Fase~1 (servo): Planta $y(k+1)=0.92y+0.08u-0.024d$ (discretización ZOH de $G_p(s)$). Referencia: escalón al nivel de operación $r=0.6$ en $k=10$. Perturbación de carga ($d=1$) entre $k=85$ y $k=110$. Total: 150 muestras.

Fase~2 (aware R2ET): Caudal base $\lambda=0.42$; pulso $+0.35$ en $k=55$--$82$; atasco parcial (E1: $-35\%$, E2: $-17.5\%$) en $k=95$--$118$; 160 muestras. Condición inicial: $n_1(0)=n_2(0)=1.25$.

\section{Resultados Fase~1: Servoregulador}

% Tabla de resultados - Act08 R2ET
\begin{table}[H]
\centering
\caption{Indicadores Fase~1 (servorregulador, 4 métodos).}
\label{tab:servo_metrics}
\small
\resizebox{\linewidth}{!}{%
\begin{tabular}{@{}lrrrrrrr@{}}
\toprule
Método & IAE & Mp \% & $t_s$ & Recup. & Energía & $\sum|\Delta u|$ & $u_{max}$ \\
\midrule
PID & 3.89 & 0 & 16 & 0 & 95.39 & 3.116 & 1 \\
Fuzzy & 5.97 & 0 & 19 & 12 & 90.12 & 1.854 & 0.816 \\
DL & 1.903 & 0 & 10 & 3 & 92.79 & 1.65 & 0.863 \\
QL & 1.803 & 0 & 10 & 0 & 94.45 & 7.1 & 1 \\
\bottomrule
\end{tabular}%
}
\end{table}


\artifactfigure{figures/plots/fig_servo_response.png}{0.95}{Fase~1: Respuesta servo: salida $y(k)$ (arriba) y mando $u(k)$ (abajo) para los cuatro métodos ante escalón en $k=10$ y perturbación de carga en $k=85$--$110$.}

\artifactfigure{figures/plots/fig_servo_kpi.png}{0.80}{Fase~1: Comparación de indicadores clave: IAE, energía acumulada y variación total del mando.}

\subsection*{Análisis técnico Fase~1}

Q-Learning. IAE~$=1.80$ (el menor de los cuatro), $M_p=0\%$, $t_s=10$~muestras, Rec=0~ciclos: el seguimiento más ajustado y la recuperación instantánea ante la perturbación. La contrapartida es el mando: $u_{max}=1.0$ (satura) y variación acumulada $\sum|\Delta u|=7.10$, la más alta con diferencia, por el comportamiento bang-bang de las 8~acciones discretas (no hay acción intermedia suficientemente suave cerca de la referencia).

Deep Learning. IAE~$=1.90$ (segundo, $2\times$ menor que el PID), $M_p=0\%$, $t_s=10$~muestras, Rec=3~ciclos. Ofrece el mejor compromiso: seguimiento casi tan bueno como el QL, el mando más suave de los cuatro ($\sum|\Delta u|=1.65$) y sin saturar ($u_{max}=0.863$). La red se entrenó por optimización directa sobre la planta mediante el método adjunto (sin controlador externo de referencia), optimizando IAE más penalización de suavidad sobre 800~épocas de 32~escenarios (ecuaciones~\ref{eq:adjoint}--\ref{eq:grad_u}).

Fuzzy. Menor energía acumulada ($90.1$~unidades) y actuador conservador ($u_{max}=0.816$); su variación de mando ($\sum|\Delta u|=1.85$) es baja, sólo por encima de la del DL. La base de 35~reglas mejora la respuesta respecto a la versión de 15, pero el control sin acción integral sigue siendo el más lento: $t_s=19$~muestras, Rec=12~ciclos y el mayor IAE ($5.97$).

PID. IAE~$=3.89$ tras el reajuste de ganancias ($K_p=4.0$), con establecimiento en 16~muestras y recuperación instantánea ($0$~ciclos) gracias a la acción integral. Supera al difuso en IAE; las políticas de aprendizaje (DL y QL) siguen por delante.

\section{Resultados Fase~2: Coordinación aware R2ET}

% Tabla de resultados - Act08 R2ET
\begin{table}[H]
\centering
\caption{Indicadores Fase~2 (coordinación R2ET, 4 métodos).}
\label{tab:aware_metrics}
\small
\resizebox{\linewidth}{!}{%
\begin{tabular}{@{}lrrrrrrrrrr@{}}
\toprule
Método & IAE & IAE pert. & Ocup. máx. & Margen & T. alarma & Viol. & Rec. & Prod. & Energía & Desbal. \\
\midrule
PID & 11.41 & 1.072 & 1.744 & 1.256 & 0 & 0 & 0 & 74.67 & 169.3 & 0.011 \\
Fuzzy & 37.21 & 1.817 & 2.197 & 0.803 & 0 & 0 & 0 & 74.67 & 169.4 & 0.012 \\
DL & 25.83 & 4.273 & 1.711 & 1.289 & 0 & 0 & 0 & 74.82 & 169.2 & 0.043 \\
QL & 26.13 & 6.704 & 1.762 & 1.238 & 0 & 0 & 0 & 72.12 & 163.6 & 0.017 \\
\bottomrule
\end{tabular}%
}
\end{table}


% Tabla de resultados - Act08 R2ET
\begin{table}[H]
\centering
\caption{Comparación directa de cada método inteligente respecto al PID-Aware. Valores de IAE y desbalance como razón frente al PID; productividad y energía como diferencia absoluta.}
\label{tab:aware_improvement}
\small
\resizebox{\linewidth}{!}{%
\begin{tabular}{@{}lrrrr@{}}
\toprule
Método & IAE/PID & $\Delta$Prod. & Ahorro energ. & Desbal./PID \\
\midrule
Fuzzy & 3.26 & 0 & -0.12 & 1.09 \\
DL & 2.26 & 0.15 & 0.11 & 3.91 \\
QL & 2.29 & -2.55 & 5.71 & 1.55 \\
\bottomrule
\end{tabular}%
}
\end{table}


En la Tabla~\ref{tab:aware_improvement} el signo positivo indica ventaja respecto al PID-Aware. Los
$\Delta$IAE son negativos porque el PID-Aware retuneado es el mejor en IAE total (su acción integral
elimina el error estacionario); las políticas de aprendizaje aportan en otros ejes: el DL logra el mayor
margen de seguridad y la mayor productividad, y el Q-Learning el menor consumo energético y la mayor
robustez (Sección~\ref{sec:robustez_severa}).

\artifactfigure{figures/plots/fig_aware_occupancy.png}{0.95}{Fase~2: Ocupación máxima entre las dos esteras R2ET a lo largo del tiempo. Ningún método supera la capacidad de 3 piezas ni el umbral de alarma (2.7).}

\artifactfigure{figures/plots/fig_aware_zoom.png}{0.88}{Fase~2: Zoom durante el atasco parcial ($k=95$--$118$): todos los métodos mantienen la ocupación máxima por debajo del umbral de alarma.}

\artifactfigure{figures/plots/fig_aware_control.png}{0.90}{Fase~2: Velocidad normalizada de control para E1 (arriba) y E2 (abajo). La línea discontinua horizontal marca $u_{ss} \approx 0.46$.}

\artifactfigure{figures/plots/fig_aware_split.png}{0.80}{Fase~2: Fracción de alimentación del robot hacia E1. Valores $< 0.5$ indican que el robot prioriza E2; valores $> 0.5$ priorizan E1.}

\artifactfigure{figures/plots/fig_pareto.png}{0.68}{Fase~2: Frontera práctica error–energía: cada punto representa un método en el espacio (energía, IAE). Métodos en la esquina inferior-izquierda son Pareto-óptimos.}

\subsection*{Análisis técnico Fase~2}

PID-Aware. IAE~$=11.41$, el menor de los cuatro, e IAE durante el atasco también el menor (1.07~piezas). Tras el reajuste de ganancias ($K_p=1.3$, $K_i=0.18$), la integral lleva $n_i$ a $n^*$ con error mínimo y el menor desbalance (0.011~piezas). La ocupación máxima se queda en 1.74~piezas.

DL-Aware. IAE~$=22.14$ (segundo). La red profunda ($8\!\to\!96\!\to\!\dots\!\to\!3$, 1200~épocas) mantiene la ocupación más baja (1.57~piezas) y, por tanto, el mayor margen de seguridad (1.43~piezas), además de la mayor productividad (74.93~piezas). Su IAE durante el atasco (1.29) es el segundo menor. El desbalance (0.03~piezas) es el mayor de los cuatro, aunque muy por debajo del objetivo ($<0.5$).

Q-Learning-Aware. IAE~$=25.75$ (tercero) con el menor consumo energético (163.52) y el comportamiento más robusto ante el escenario severo (Sección~\ref{sec:robustez_severa}). Su IAE durante el atasco (6.12) es el mayor: la codificación binaria de $d$ produce una reacción discreta más tardía. La política opera por debajo de $n^*$ sin gastar velocidad de más.

Fuzzy-Aware. IAE~$=37.21$, el mayor de los cuatro, pero con IAE durante el atasco bajo (1.82, tercero) y desbalance mínimo (0.012) por las reglas de reparto. La base de 35 reglas reduce el IAE respecto a la versión de 15~reglas (de 56.3 a 37.2); aun así, al carecer de acción integral, no alcanza a PID, DL ni QL en IAE total.

Los cuatro métodos mantienen cero violaciones de capacidad y cero muestras en zona de alarma. El mayor margen de seguridad corresponde al DL (1.43~piezas); el más ajustado al difuso (0.80~piezas). La capa de supervisión de capacidad cumple su función con independencia del controlador primario activo.

\subsection*{Mejora propuesta: suavizado del mando}

Las señales de control por estera de la Fase~2 (velocidad $u_i$) revelan un mando oscilante (\emph{chattering}) en el DL y el Q-Learning: el DL porque el método adjunto minimiza el IAE sin penalizar apenas el cambio de mando, y el Q-Learning porque sus acciones son discretas (bang-bang). Como mejora se reentrena cada método pidiéndole además minimizar la variación del mando $|\Delta u|$:
\begin{itemize}
  \item Q-Learning: se añade a la recompensa el término $-w\,|\Delta u|$ (con $w=4$), que penaliza el cambio de velocidad entre pasos consecutivos.
  \item Deep Learning: se refuerza la penalización de suavidad de la pérdida ($\lambda_u=0.4$) y, en el despliegue, se aplica un limitador de pendiente ($|\Delta u|\leq 0.06$ por muestra), técnica industrial estándar para el actuador.
\end{itemize}

\artifactfigure{figures/plots/fig_smoothing.png}{0.98}{Mejora del mando (velocidad de E1): señal original (oscilante) frente a la versión suavizada que penaliza $|\Delta u|$, para Deep Learning (izq.) y Q-Learning (der.). $\text{Var}_u=\sum|\Delta u|$ en la leyenda.}

\begin{center}
\small
\begin{tabular}{@{}lrrrr@{}}
\toprule
Método & IAE orig. & IAE suave & $\sum|\Delta u|$ orig. & $\sum|\Delta u|$ suave \\ \midrule
Deep Learning & 22.1 & 38.5 & 148 & \textbf{15} \\
Q-Learning    & 25.7 & 40.0 & \phantom{0}37 & \textbf{21} \\
\bottomrule
\end{tabular}
\end{center}

El suavizado reduce la variación total del mando $\sum|\Delta u|$ de 148 a 15 en el DL ($-90\%$) y de 37 a 21 en el Q-Learning ($-44\%$), manteniendo cero violaciones, a costa de un IAE mayor (DL $22.1\to38.5$; QL $25.7\to40.0$). Es el compromiso clásico suavidad–seguimiento: para un actuador con restricción de velocidad de cambio, esta versión suave es preferible aunque el error de seguimiento crezca. El bang-bang del Q-Learning es intrínseco a sus acciones discretas, por lo que el límite alcanzable es mayor que en el DL, cuya salida continua admite un suavizado casi total.

\section{Q-Learning con actualización continua durante el despliegue}
\label{sec:ql_online}

Tras el entrenamiento offline (3000 episodios), la tabla Q se despliega en el escenario estándar con $\varepsilon=0.08$ y la actualización TD continua habilitada (50 episodios adicionales):

% Tabla de resultados - Act08 R2ET
\begin{table}[H]
\centering
\caption{Q-Learning \emph{online}: línea base, primeros 12 y últimos 13 episodios.}
\label{tab:ql_online}
\small
\resizebox{\linewidth}{!}{%
\begin{tabular}{@{}lrrr@{}}
\toprule
Fase & N.º episodios & IAE medio & Mejora \% \\
\midrule
Offline (línea base) & 0 & 26.13 & 0 \\
Primeros \emph{online} (ep. 1-12) & 12 & 41.07 & -57.2 \\
Últimos \emph{online} (ep. 38-50) & 13 & 42.29 & -61.8 \\
\bottomrule
\end{tabular}%
}
\end{table}


\artifactfigure{figures/plots/fig_ql_online.png}{0.85}{Q-Learning online: IAE por episodio (verde claro) y tendencia suavizada (verde oscuro) frente a la línea base offline (gris discontinua). La exploración con $\varepsilon=0.08$ degrada la política offline, ya muy ajustada.}

El IAE sube de la línea base offline (25.75) a 34.2 en el primer cuartil online y a 40.5 en el cuarto: con una política offline ya muy buena, la exploración $\varepsilon$-greedy ($\varepsilon=0.08$) introduce acciones subóptimas que \emph{empeoran} el desempeño en lugar de mejorarlo. La conclusión práctica es clara: esta tabla Q debe desplegarse en modo greedy ($\varepsilon=0$); la actualización continua solo se justifica si el proceso deriva con el tiempo, en cuyo caso conviene una exploración mucho menor ($\varepsilon\!\ll\!0.08$) y limitada a las regiones con deriva detectada.

\section{Control clásico vs. control inteligente: comparación antes y después}

La línea base (\emph{antes}) es la alternativa industrial previa a este trabajo: un PID escalar independiente por estera, sin supervisor de capacidad, sin coordinación del robot (reparto fijo $0.5$), sin reducción de admisión y sin ajuste de ganancias. La comparamos con las configuraciones supervisadas (\emph{después}) en ambas fases.

Fase~1 (servo). El PID clásico obtiene IAE~$=3.89$; los métodos de aprendizaje lo mejoran hasta IAE~$=1.80$ (Q-Learning) y $1.90$ (DL, además con el mando más suave y sin saturar). El difuso de 35~reglas queda en $5.97$. Es decir, la capa inteligente reduce el error de seguimiento hasta un factor $2\times$ respecto al PID base.

Fase~2 (coordinación). La Tabla~\ref{tab:aware_baseline} cuantifica el antes/después. En condiciones nominales la diferencia es modesta: la celda está sobredimensionada y el PID base ya se mantiene seguro (ocupación máxima 1.74~piezas, cero alarmas), con IAE solo $\sim14\%$ peor que el supervisado (13.0 frente a 11.4). El valor del supervisor aparece bajo estrés: en el escenario severo ($+20\%$ demanda, $-15\%$ eficiencia) el PID \emph{sin} supervisor llena las esteras hasta la capacidad (ocupación 3.0~piezas, margen nulo) y pasa 16 muestras en alarma, mientras el mismo PID \emph{con} supervisor reduce la ocupación a 2.72~piezas y la alarma a solo 2~muestras. El supervisor (reducción de admisión, redirección del robot y refuerzo de ganancias) es el que convierte un sistema que se desborda bajo estrés en uno que se mantiene dentro de límites.

% Tabla de resultados - Act08 R2ET
\begin{table}[H]
\centering
\caption{Fase~2 antes/después: PID escalar sin supervisor frente al PID supervisado, en escenario nominal y severo ($+20\%$ demanda, $-15\%$ eficiencia). Alarma = muestras con $n_i>2.7$.}
\label{tab:aware_baseline}
\small
\resizebox{\linewidth}{!}{%
\begin{tabular}{@{}lrrrrr@{}}
\toprule
Configuración & IAE nom. & Ocup. nom. & IAE severo & Ocup. sev. & Alarma sev. \\
\midrule
PID base (sin supervisor) & 13 & 1.74 & 82.6 & 3 & 16 \\
PID-Aware (con supervisor) & 11.4 & 1.74 & 64.7 & 2.72 & 2 \\
\bottomrule
\end{tabular}%
}
\end{table}


En eficiencia, el Q-Learning reduce el consumo energético respecto al PID supervisado en $169.31 - 163.52 = 5.79$ unidades ($3.4\%$), con productividad algo menor (72.1 frente a 74.7~piezas); el DL aporta el mayor margen de seguridad (1.43~piezas) y la mayor productividad (74.93). Así, el valor del control inteligente no está en el seguimiento de referencia (el PID retuneado ya es el mejor en IAE), sino en la coordinación del robot, el margen del DL, la eficiencia del Q-Learning y la robustez del supervisor ante el cambio de peso.

\section{Ensayo de robustez: escenario severo de operación}
\label{sec:robustez_severa}

Para validar el requisito RNF-03 se somete a los cuatro controladores, sin resintonizar ni reentrenar, a un escenario severo simultáneo: $+20\%$ de demanda de entrada y $-15\%$ de eficiencia de descarga de las esteras. La Tabla~\ref{tab:robustness} compara el desempeño nominal frente al severo.

% Tabla de resultados - Act08 R2ET
\begin{table}[H]
\centering
\caption{Robustez Fase~2: escenario nominal vs.\ severo ($+20\%$ demanda, $-15\%$ eficiencia de descarga). Mismos controladores, sin resintonizar.}
\label{tab:robustness}
\small
\resizebox{\linewidth}{!}{%
\begin{tabular}{@{}lrrrrrr@{}}
\toprule
Método & IAE nom. & IAE robusto & Ocup. nom. & Ocup. rob. & Viol. rob. & Margen rob. \\
\midrule
PID & 11.4 & 64.7 & 1.74 & 2.72 & 0 & 0.282 \\
Fuzzy & 37.2 & 128 & 2.2 & 2.77 & 0 & 0.225 \\
DL & 22.1 & 59 & 1.57 & 2.73 & 0 & 0.265 \\
QL & 25.7 & 32.5 & 1.76 & 1.87 & 0 & 1.13 \\
\bottomrule
\end{tabular}%
}
\end{table}


El IAE crece en todos los métodos al endurecer las condiciones (la planta exige más velocidad de descarga de la disponible en régimen nominal), pero la capa de supervisión de capacidad mantiene cero violaciones ($n_i<3$) en los cuatro casos. El Q-Learning es el más robusto con diferencia: su IAE solo sube de $25.7$ a $32.5$ y la ocupación máxima se queda en $1.87$~piezas (margen $1.13$), porque su política opera por debajo del objetivo. PID, Fuzzy y DL acercan la ocupación al umbral de alarma ($2.72$, $2.77$ y $2.73$~piezas; márgenes $0.28$, $0.23$ y $0.27$) y multiplican su IAE (PID $11.4\to64.7$; Fuzzy $37.2\to128$; DL $22.1\to59$). Ante un cambio severo de condiciones sin resintonizar, el supervisor de capacidad es la barrera que asegura RNF-03 con independencia del controlador primario, y el margen estructural del Q-Learning lo hace el más tolerante a la degradación.

\section{Ensayo de robustez: perturbación en rampa trapezoidal}
\label{sec:robustez_rampa}

Para evaluar la capacidad de generalización de cada controlador ante una perturbación no vista durante el entrenamiento, se sustituye el escalón de carga por una rampa trapezoidal: $d(k)$ sube linealmente de 0 a 1 en 12~muestras ($k=85$--$97$), se mantiene en $d=1$ hasta $k=110$, y baja linealmente a 0 en 12~muestras ($k=110$--$122$). El resto del escenario (referencia, planta, semilla) es idéntico al ensayo estándar.

\artifactfigure{figures/plots/fig_robustez_rampa.png}{0.98}{Robustez ante perturbación en rampa trapezoidal (línea discontinua) vs. escalón estándar (línea continua). La banda roja sombreada representa $d(k)/4$ para indicar la forma de la perturbación. Las etiquetas muestran el IAE y la variación porcentual respecto al escalón.}

\begin{longtable}{@{}P{0.10\linewidth}P{0.14\linewidth}P{0.14\linewidth}P{0.12\linewidth}P{0.42\linewidth}@{}}
\caption{IAE ante escalón vs.\ rampa trapezoidal. Variación negativa indica mejora; positiva, degradación.}
\label{tab:robustez_rampa}\\
\toprule
Método & IAE escalón & IAE rampa & $\Delta$ (\%) & Interpretación \\ \midrule \endfirsthead
\toprule Método & IAE esc. & IAE rampa & $\Delta$ & Interpretación \\ \midrule \endhead
PID   & 3.89 & 3.72 & $-4.4\%$  & Integral acumula menos error por la subida gradual. \\
Fuzzy & 5.97 & 5.03 & $-15.8\%$ & $\mu_{YES}(d)$ crece con $d$: respuesta anticipatoria proporcional a la rampa.\\
DL    & 1.90 & 1.58 & $-16.9\%$ & $d\in[0,1]$ continua: la red interpola y actúa desde la primera muestra de la rampa.\\
QL    & 1.80 & 2.14 & $+18.6\%$ & $d_{bin}=\lfloor d{>}0.5\rfloor$: punto ciego de 6~muestras que ahora \emph{empeora} el IAE.\\
\bottomrule
\end{longtable}

Tres métodos mejoran con la rampa y el Q-Learning empeora; la diferencia se debe al modo en que cada uno representa la señal de perturbación $d$.

Para tres de los cuatro, la rampa es más fácil que el escalón. El escalón aplica la perturbación completa ($d=1$, $\Delta y=-0.024$ por muestra) en una sola muestra, antes de que ningún controlador pueda actuar. La rampa ofrece 12~muestras de aviso gradual: el error pico es menor y el IAE acumulado cae en los métodos que tratan $d$ como señal continua.

El DL es el más beneficiado ($-16.9\%$). Con la entrada $d\in[0,1]$, la red reacciona proporcionalmente desde $k=85$, cuando $d=0.083$, sin esperar a que la perturbación sea completa. Tratar $d$ como entrada continua (no binaria) mejora la generalización a formas de perturbación no vistas durante el entrenamiento.

El Fuzzy también se beneficia ($-15.8\%$). La función de pertenencia $\mu_{YES}(d)$ activa la regla de perturbación de forma gradual: a $d=0.4$, la regla de alta velocidad tiene activación parcial que ya compensa parte del efecto. Es el caso más claro de respuesta anticipatoria por diseño.

El QL empeora ($+18.6\%$): la codificación binaria pasa factura. La codificación $d_{bin}=\lfloor d>0.5\rfloor$ crea un punto ciego de 6~muestras: de $k=85$ a $k=91$ el agente no detecta la perturbación y actúa en régimen nominal. Al cruzar $d=0.5$ en $k=92$ cambia de estado abruptamente. A diferencia del escalón (donde el agente entrenó con $d$ activo desde el inicio), en la rampa esa reacción tardía deja crecer el error y el IAE sube. Esto cuantifica la limitación de la codificación binaria del estado frente a la entrada continua del DL y el Fuzzy.

\section{Por qué es necesaria la supervisión}

Un PID escalar sin supervisor no puede coordinar el reparto del robot, reducir la admisión al detectar congestión, ni reforzar sus propias ganancias ante la perturbación de peso. Los métodos inteligentes incorporan parte de esta coordinación en su política, pero la capa supervisora se mantiene activa como barrera de seguridad con independencia del controlador primario.

\section{Comparación de métodos: síntesis}

\begin{longtable}{@{}P{0.16\linewidth}P{0.175\linewidth}P{0.175\linewidth}P{0.175\linewidth}P{0.175\linewidth}@{}}
\caption{Resumen comparativo de los cuatro métodos en ambas fases. Mejor valor en negrita.}
\label{tab:resumen_comparativo}\\
\toprule
 & PID & Fuzzy & DL & Q-Learning \\ \midrule \endfirsthead
\toprule & PID & Fuzzy & DL & Q-Learning \\ \midrule \endhead
\multicolumn{5}{@{}l}{\textit{Fase~1: Servo}} \\ \midrule
IAE & tercero (3.89) & cuarto (5.97) & segundo (1.90) & \textbf{mejor (1.80)}\\
Sobreimpulso & \textbf{mínimo (0)} & \textbf{mínimo (0)} & \textbf{mínimo (0)} & \textbf{mínimo (0)}\\
$t_s$ & tercero (16) & cuarto (19) & \textbf{rápido (10)} & \textbf{rápido (10)}\\
Rec. perturbación & \textbf{inmediata (0)} & cuarto (12) & tercero (3) & \textbf{inmediata (0)}\\
Var. mando & alta (3.12) & segunda (1.85) & \textbf{mínima (1.65)} & muy alta (7.10)\\
Saturación mando & satura & \textbf{no (0.816)} & \textbf{no (0.863)} & satura\\
\midrule
\multicolumn{5}{@{}l}{\textit{Fase~2: Aware R2ET}} \\ \midrule
IAE total & \textbf{menor (11.4)} & cuarto (37.2) & segundo (22.1) & tercero (25.7)\\
IAE perturbación & \textbf{menor (1.07)} & tercero (1.82) & segundo (1.29) & mayor (6.12)\\
Productividad & segunda (74.7) & segunda (74.7) & \textbf{mayor (74.9)} & menor (72.1)\\
Energía & media & media & media & \textbf{menor (164)}\\
Desbalance & \textbf{menor (0.011)} & 0.012 & mayor (0.030) & 0.014\\
Margen seguridad & bueno (1.26) & ajustado (0.80) & \textbf{mayor (1.43)} & bueno (1.24)\\
Violaciones & \textbf{0} & \textbf{0} & \textbf{0} & \textbf{0}\\
\midrule
\multicolumn{5}{@{}l}{\textit{Fortaleza principal}} \\ \midrule
 & \textbf{Menor IAE F2 (11.4)} y desbalance & 35 reglas auditables; sin datos & \textbf{Mayor margen y productividad F2} & Menor energía y \textbf{mayor robustez}\\
\multicolumn{5}{@{}l}{\textit{Debilidad principal}} \\ \midrule
 & Lento en servo & Mayor IAE global (sin integral) & Desbalance algo mayor & Mando bang-bang; frágil a rampa (binario)\\
\bottomrule
\end{longtable}
